% Topic 16: Elements inside HTML Table
\section{Table-এর ভিতরে List, Image, Hyperlink ও Nested Table}

\begin{learningbox}
এই টপিক শেষে শিক্ষার্থী পারবে:
\begin{itemize}
\item HTML table cell-এর ভিতরে list, image, hyperlink ও formatting tag ব্যবহার করতে।
\item table cell-এর ভিতরে ordered list ও unordered list বসাতে।
\item image hyperlink তৈরি করতে।
\item Nested Table কী তা ব্যাখ্যা করতে এবং basic nested table code লিখতে।
\item HSC CQ/practical question-এ table-এর ভিতরে multiple HTML element ব্যবহার করে code লিখতে।
\end{itemize}
\end{learningbox}

HTML table-এর cell-এর ভিতরে শুধু text রাখা যায় না; সেখানে list, image, hyperlink, formatting tag, এমনকি আরেকটি table-ও রাখা যায়। তাই HSC ICT Chapter 4-এর practical/CQ অংশে table-এর ভিতরে list, image, link, underline text, rowspan/colspan--এসব একসাথে ব্যবহার করে code লিখতে বলা হয়।

\subsection{Table Cell-এর ভিতরে কী কী রাখা যায়?}

HTML table-এর \texttt{<td>} বা \texttt{<th>} cell-এর ভিতরে অনেক ধরনের HTML element রাখা যায়।

\begin{center}
\footnotesize
\begin{tabular}{|p{0.30\textwidth}|p{0.52\textwidth}|}
\hline
\textbf{Cell-এর ভিতরে রাখা যায়} & \textbf{Example} \\
\hline
সাধারণ text & \texttt{<td>ICT</td>} \\
\hline
Bold/underline text & \texttt{<td><u>ICT</u></td>} \\
\hline
Hyperlink & \texttt{<td><a href="board.html">Board</a></td>} \\
\hline
Image & \texttt{<td><img src="book.jpg" width="100"></td>} \\
\hline
Ordered list & \texttt{<td><ol><li>Book</li></ol></td>} \\
\hline
Unordered list & \texttt{<td><ul><li>Pen</li></ul></td>} \\
\hline
Nested table & \texttt{<td><table>...</table></td>} \\
\hline
\end{tabular}
\end{center}

\begin{keypointbox}{মূল ধারণা}
Table cell-এর ভিতরে content রাখার সময় content অবশ্যই \texttt{<td>} বা \texttt{<th>} tag-এর ভিতরে রাখতে হবে। \texttt{<table>} tag-এর ভিতরে সরাসরি \texttt{<ol>}, \texttt{<ul>}, \texttt{<img>} বা \texttt{<a>} লেখা ঠিক নয়।
\end{keypointbox}

\subsection{Table-এর ভিতরে List}

কোনো table cell-এর ভিতরে ordered list বা unordered list বসাতে হলে \texttt{<td>} tag-এর ভিতরে \texttt{<ol>} বা \texttt{<ul>} লিখতে হয়।

\subsubsection{সাধারণ Syntax}

\begin{examplebox}{Code}
\begin{verbatim}
<td>
  <ol>
    <li>Item 1</li>
    <li>Item 2</li>
  </ol>
</td>
\end{verbatim}
\end{examplebox}

অথবা--

\begin{examplebox}{Code}
\begin{verbatim}
<td>
  <ul>
    <li>Item 1</li>
    <li>Item 2</li>
  </ul>
</td>
\end{verbatim}
\end{examplebox}

\subsection{Example: Ordered List ও Unordered List Table}

\begin{examplebox}{Output ধারণা}
\begin{center}
\begin{tabular}{|p{0.36\textwidth}|p{0.36\textwidth}|}
\hline
\multicolumn{2}{|c|}{\textbf{A \& H Company Ltd.}} \\
\hline
\textbf{Order List} & \textbf{Unorder List} \\
\hline
i. Book \newline ii. Paper \newline iii. Note Book & Marker \newline Ink \newline Pencil \\
\hline
\end{tabular}
\end{center}
\end{examplebox}

\begin{examplebox}{Code}
\begin{verbatim}
<!DOCTYPE html>
<html>
<head>
  <title>HTML List in Table</title>
</head>
<body>
  <table border="2" cellspacing="0" cellpadding="5">
    <tr>
      <th colspan="2" align="center">A & H Company Ltd.</th>
    </tr>
    <tr>
      <th>Order List</th>
      <th>Unorder List</th>
    </tr>
    <tr>
      <td>
        <ol type="i">
          <li>Book</li>
          <li>Paper</li>
          <li>Note Book</li>
        </ol>
      </td>
      <td>
        <ul type="square">
          <li>Marker</li>
          <li>Ink</li>
          <li>Pencil</li>
        </ul>
      </td>
    </tr>
  </table>
</body>
</html>
\end{verbatim}
\end{examplebox}

\begin{boardbox}{Exam note}
HSC board-style CQ practice-এ table-এর cell-এর ভিতরে ordered list ও unordered list বসিয়ে code লিখতে বলা হয়। তাই \texttt{<td>} tag-এর ভিতরে \texttt{<ol>} বা \texttt{<ul>} বসানোর নিয়ম ভালোভাবে practice করতে হবে।
\end{boardbox}

\subsection{Table-এর ভিতরে Hyperlink}

Table cell-এর ভিতরে link বসাতে \texttt{<td>} tag-এর ভিতরে \texttt{<a href="">...</a>} লিখতে হয়।

\subsubsection{Syntax}

\begin{examplebox}{Code}
\begin{verbatim}
<td>
  <a href="url">Link Text</a>
</td>
\end{verbatim}
\end{examplebox}

\subsubsection{Example}

\begin{examplebox}{Code}
\begin{verbatim}
<td>
  <a href="https://www.educationboard.gov.bd">Board</a>
</td>
\end{verbatim}
\end{examplebox}

নতুন tab/window-এ link open করতে \texttt{target="\_blank"} attribute ব্যবহার করা হয়।

\subsection{Example: Table Cell-এর ভিতরে Link}

\begin{examplebox}{Output ধারণা}
\begin{center}
\begin{tabular}{|c|c|}
\hline
\textbf{Subject} & \textbf{Board} \\
\hline
ICT & Board \\
\hline
\end{tabular}
\end{center}
Board text-এ click করলে নির্দিষ্ট website open হবে।
\end{examplebox}

\begin{examplebox}{Code}
\begin{verbatim}
<!DOCTYPE html>
<html>
<head>
  <title>Link in Table</title>
</head>
<body>
  <table border="1" cellspacing="0" cellpadding="5">
    <tr>
      <th>Subject</th>
      <th>Board</th>
    </tr>
    <tr>
      <td>ICT</td>
      <td>
        <a href="https://www.educationboard.gov.bd" target="_blank">
          Board
        </a>
      </td>
    </tr>
  </table>
</body>
</html>
\end{verbatim}
\end{examplebox}

\subsection{Table-এর ভিতরে Image}

Table cell-এর ভিতরে image বসাতে \texttt{<td>} tag-এর ভিতরে \texttt{<img>} tag লিখতে হয়।

\subsubsection{Syntax}

\begin{examplebox}{Code}
\begin{verbatim}
<td>
  <img src="image.jpg" alt="Image" width="100" height="80">
</td>
\end{verbatim}
\end{examplebox}

এখানে \texttt{src} দ্বারা ছবির source নির্ধারণ করা হয়, \texttt{alt} দ্বারা বিকল্প text দেওয়া হয় এবং \texttt{width}, \texttt{height} দ্বারা ছবির প্রস্থ ও উচ্চতা নির্ধারণ করা হয়।

\subsection{Example: Book Image in Table}

\begin{examplebox}{Output ধারণা}
\begin{center}
\begin{tabular}{|c|c|}
\hline
\textbf{Book} & \textbf{Image} \\
\hline
ICT Book & book.jpg \\
\hline
\end{tabular}
\end{center}
\end{examplebox}

\begin{examplebox}{Code}
\begin{verbatim}
<!DOCTYPE html>
<html>
<head>
  <title>Image in Table</title>
</head>
<body>
  <table border="1" cellspacing="0" cellpadding="5">
    <tr>
      <th>Book</th>
      <th>Image</th>
    </tr>
    <tr>
      <td>ICT Book</td>
      <td>
        <img src="book.jpg" alt="ICT Book" width="100" height="80">
      </td>
    </tr>
  </table>
</body>
</html>
\end{verbatim}
\end{examplebox}

\subsection{Table-এর ভিতরে Image Hyperlink}

যদি table cell-এর ভিতরের image-এ click করলে কোনো page বা website open হয়, তাহলে \texttt{<a>} tag-এর ভিতরে \texttt{<img>} tag বসাতে হবে।

\subsubsection{Syntax}

\begin{examplebox}{Code}
\begin{verbatim}
<td>
  <a href="url">
    <img src="image.jpg" alt="Image">
  </a>
</td>
\end{verbatim}
\end{examplebox}

\subsection{Example: Book Image-এ Click করলে Website Open}

\begin{examplebox}{Code}
\begin{verbatim}
<!DOCTYPE html>
<html>
<head>
  <title>Image Link in Table</title>
</head>
<body>
  <table border="1" cellspacing="0" cellpadding="5">
    <tr>
      <th>Book</th>
      <th>Image Link</th>
    </tr>
    <tr>
      <td>ICT Book</td>
      <td>
        <a href="https://www.book.com" target="_blank">
          <img src="book.jpg" alt="ICT Book" width="100" height="80">
        </a>
      </td>
    </tr>
  </table>
</body>
</html>
\end{verbatim}
\end{examplebox}

\begin{keypointbox}{মনে রাখবে}
Clickable text তৈরি করতে \texttt{<a>} tag-এর ভিতরে text লিখতে হয়। Clickable image তৈরি করতে \texttt{<a>} tag-এর ভিতরে \texttt{<img>} tag লিখতে হয়।
\end{keypointbox}

\subsection{Table-এর ভিতরে List + Hyperlink একসাথে}

List item-এর ভিতরেও link বসানো যায়। যেমন Fruits item-এ click করলে কোনো website নতুন tab-এ open হবে।

\begin{examplebox}{Code}
\begin{verbatim}
<!DOCTYPE html>
<html>
<head>
  <title>List and Link</title>
</head>
<body>
  <ul type="circle">
    <li>Vegetables</li>
    <li>
      <a href="https://www.fruitbox.com" target="_blank">Fruits</a>
      <ol type="i" start="4">
        <li>Ripe</li>
        <li>Dried</li>
      </ol>
    </li>
  </ul>
</body>
</html>
\end{verbatim}
\end{examplebox}

\subsection{Nested Table কী?}

\begin{definitionbox}{Nested Table}
HTML document-এ একটি table-এর \texttt{<td>} বা \texttt{<th>} cell-এর ভিতরে আরেকটি table তৈরি করা হলে তাকে Nested Table বলে।
\end{definitionbox}

Nested table technically করা যায়। তবে বেশি nested table ব্যবহার করলে code জটিল হয় এবং accessibility ও responsive design-এ সমস্যা হতে পারে। তাই প্রয়োজন ছাড়া nested table ব্যবহার করা উচিত নয়।

\subsection{Nested Table-এর Syntax}

\begin{examplebox}{Code}
\begin{verbatim}
<table border="1">
  <tr>
    <td>
      Outer Table Cell
    </td>
    <td>
      <table border="1">
        <tr>
          <td>Inner Table Cell</td>
        </tr>
      </table>
    </td>
  </tr>
</table>
\end{verbatim}
\end{examplebox}

\subsection{Example: Student Info with Nested Marks Table}

\begin{examplebox}{Output ধারণা}
\begin{center}
\begin{tabular}{|c|c|}
\hline
\textbf{Student} & \textbf{Marks} \\
\hline
Rafi & Bangla 85 \newline English 90 \\
\hline
\end{tabular}
\end{center}
\end{examplebox}

\begin{examplebox}{Code}
\begin{verbatim}
<!DOCTYPE html>
<html>
<head>
  <title>Nested Table Example</title>
</head>
<body>
  <table border="1" cellspacing="0" cellpadding="5">
    <tr>
      <th>Student</th>
      <th>Marks</th>
    </tr>
    <tr>
      <td>Rafi</td>
      <td>
        <table border="1" cellspacing="0" cellpadding="5">
          <tr>
            <th>Subject</th>
            <th>Mark</th>
          </tr>
          <tr>
            <td>Bangla</td>
            <td>85</td>
          </tr>
          <tr>
            <td>English</td>
            <td>90</td>
          </tr>
        </table>
      </td>
    </tr>
  </table>
</body>
</html>
\end{verbatim}
\end{examplebox}

\subsection{Nested Table ব্যবহারে সতর্কতা}

\begin{mistakebox}{সতর্কতা}
\begin{itemize}
\item Nested table code পড়া কঠিন হয়।
\item ভুল closing tag দিলে layout নষ্ট হয়।
\item accessibility সমস্যা হতে পারে।
\item mobile responsive design-এ সমস্যা হতে পারে।
\item শুধুমাত্র tabular data-এর ভিতরে sub-table দরকার হলে ব্যবহার করা উচিত।
\end{itemize}
\end{mistakebox}

\begin{boardbox}{Book note}
HSC exam-এর জন্য nested table জানা ভালো, কিন্তু professional web design-এ layout তৈরির জন্য table ব্যবহার না করে CSS layout ব্যবহার করা উত্তম।
\end{boardbox}

\subsection{Table-এর ভিতরে Multiple Elements}

একটি cell-এর ভিতরে একাধিক element একসাথে থাকতে পারে।

\begin{examplebox}{Code}
\begin{verbatim}
<td>
  <h3>ICT Book</h3>
  <img src="book.jpg" alt="ICT Book" width="100">
  <p>Price: 300 Tk</p>
  <a href="book.html">Details</a>
</td>
\end{verbatim}
\end{examplebox}

এখানে একই cell-এর ভিতরে heading, image, paragraph ও link আছে।

\subsection{Full Practice Example: Table + List + Image + Link}

\begin{examplebox}{Output ধারণা}
\begin{center}
\begin{tabular}{|c|c|}
\hline
\textbf{Item} & \textbf{Details} \\
\hline
Book & image + link \\
\hline
Subjects & ordered list \\
\hline
\end{tabular}
\end{center}
\end{examplebox}

\begin{examplebox}{Code}
\begin{verbatim}
<!DOCTYPE html>
<html>
<head>
  <title>Table with Multiple Elements</title>
</head>
<body>
  <table border="1" cellspacing="0" cellpadding="5">
    <caption>ICT Learning Materials</caption>
    <tr>
      <th>Item</th>
      <th>Details</th>
    </tr>
    <tr>
      <td>Book</td>
      <td>
        <a href="https://www.book.com" target="_blank">
          <img src="book.jpg" alt="ICT Book" width="100" height="80">
        </a>
        <br>
        <a href="https://www.book.com" target="_blank">Buy Book</a>
      </td>
    </tr>
    <tr>
      <td>Subjects</td>
      <td>
        <ol type="1">
          <li>Web Design</li>
          <li>HTML</li>
          <li>Database</li>
        </ol>
      </td>
    </tr>
  </table>
</body>
</html>
\end{verbatim}
\end{examplebox}

\subsection{Board/Exam Focus}

\begin{boardbox}{বেশি গুরুত্বপূর্ণ}
\begin{itemize}
\item Table-এর cell-এর ভিতরে ordered/unordered list বসানো।
\item Table-এর cell-এর ভিতরে hyperlink তৈরি করা।
\item Table-এর cell-এর ভিতরে image বসানো।
\item Image-কে link হিসেবে ব্যবহার করা।
\item Table-এর cell-এর ভিতরে underline/bold formatting ব্যবহার করা।
\item \texttt{rowspan}, \texttt{colspan}, \texttt{align}, \texttt{target}, \texttt{src}, \texttt{width}, \texttt{height} একসাথে ব্যবহার করা।
\item Figure দেখে HTML code লেখা।
\end{itemize}
\end{boardbox}

\subsection{Exam-style Short Answer}

\begin{enumerate}
\item \textbf{Table-এর cell-এর ভিতরে list কীভাবে লেখা যায়?}\\
উত্তর: Table-এর \texttt{<td>} বা \texttt{<th>} tag-এর ভিতরে \texttt{<ol>} অথবা \texttt{<ul>} tag ব্যবহার করে list লেখা যায়। প্রতিটি item লেখার জন্য \texttt{<li>} tag ব্যবহার করতে হয়।

\item \textbf{Table-এর ভিতরে hyperlink কীভাবে তৈরি করা যায়?}\\
উত্তর: Table cell-এর ভিতরে \texttt{<a href="url">Link Text</a>} লিখে hyperlink তৈরি করা যায়। যেমন: \texttt{<td><a href="board.html">Board</a></td>}।

\item \textbf{Table-এর ভিতরে image কীভাবে বসানো যায়?}\\
উত্তর: Table cell-এর ভিতরে \texttt{<img>} tag ব্যবহার করে image বসানো যায়। যেমন: \texttt{<td><img src="book.jpg" alt="Book" width="100"></td>}।

\item \textbf{Nested Table কী?}\\
উত্তর: একটি table-এর \texttt{<td>} বা \texttt{<th>} cell-এর ভিতরে আরেকটি সম্পূর্ণ table তৈরি করা হলে তাকে Nested Table বলে।

\item \textbf{Image hyperlink কীভাবে তৈরি করা যায়?}\\
উত্তর: Image hyperlink তৈরি করতে \texttt{<a>} tag-এর ভিতরে \texttt{<img>} tag বসাতে হয়। যেমন: \texttt{<a href="index.html"><img src="logo.png" alt="Logo"></a>}।
\end{enumerate}

\subsection{Practice MCQ}

\begin{enumerate}
\item Table cell-এর ভিতরে image বসাতে কোন tag ব্যবহার হয়?\\
ক) \texttt{<a>} \quad খ) \texttt{<img>} \quad গ) \texttt{<ol>} \quad ঘ) \texttt{<tr>}

\item Table cell-এর ভিতরে link তৈরি করতে কোন tag ব্যবহার হয়?\\
ক) \texttt{<a>} \quad খ) \texttt{<td>} \quad গ) \texttt{<br>} \quad ঘ) \texttt{<hr>}

\item Table cell-এর ভিতরে ordered list বসাতে কোন tag ব্যবহার হবে?\\
ক) \texttt{<ul>} \quad খ) \texttt{<ol>} \quad গ) \texttt{<img>} \quad ঘ) \texttt{<caption>}

\item Image hyperlink তৈরির সঠিক structure কোনটি?\\
ক) \texttt{<img href="home.html">}\\
খ) \texttt{<a href="home.html"><img src="logo.png" alt="Logo"></a>}\\
গ) \texttt{<a src="logo.png">Home</a>}\\
ঘ) \texttt{<table img="logo.png">}

\item একটি table-এর cell-এর ভিতরে আরেকটি table থাকলে তাকে কী বলে?\\
ক) Empty Table \quad খ) Nested Table \quad গ) Ordered Table \quad ঘ) External Table
\end{enumerate}

\begin{boardbox}{MCQ Answer Key}
১-খ, ২-ক, ৩-খ, ৪-খ, ৫-খ
\end{boardbox}

\subsection{CQ Practice}

\begin{practicebox}
\textbf{উদ্দীপক:} একটি webpage-এ ``ICT Learning Materials'' নামে একটি table তৈরি করতে হবে। Table-এ দুইটি column থাকবে: Item ও Details। Book row-এর Details cell-এ \texttt{book.jpg} image থাকবে এবং image-এ click করলে \texttt{book.html} page open হবে। Subjects row-এর Details cell-এ ordered list হিসেবে HTML, CSS, JavaScript থাকবে।

\textbf{প্রশ্ন}\\
ক. Nested Table কী?\\
খ. Image hyperlink তৈরির নিয়ম ব্যাখ্যা কর।\\
গ. উদ্দীপকের জন্য HTML code লেখ।\\
ঘ. Table-এর cell-এর ভিতরে list, image ও hyperlink ব্যবহারের সুবিধা বিশ্লেষণ কর।
\end{practicebox}

\subsubsection{গ. Code}

\begin{examplebox}{Code}
\begin{verbatim}
<!DOCTYPE html>
<html>
<head>
  <title>ICT Learning Materials</title>
</head>
<body>
  <table border="1" cellspacing="0" cellpadding="5">
    <caption>ICT Learning Materials</caption>
    <tr>
      <th>Item</th>
      <th>Details</th>
    </tr>
    <tr>
      <td>Book</td>
      <td>
        <a href="book.html">
          <img src="book.jpg" alt="ICT Book" width="100" height="80">
        </a>
      </td>
    </tr>
    <tr>
      <td>Subjects</td>
      <td>
        <ol>
          <li>HTML</li>
          <li>CSS</li>
          <li>JavaScript</li>
        </ol>
      </td>
    </tr>
  </table>
</body>
</html>
\end{verbatim}
\end{examplebox}

\subsubsection{ঘ. বিশ্লেষণ}

Table-এর cell-এর ভিতরে list ব্যবহার করলে তথ্য category-wise সুন্দরভাবে সাজানো যায়। Image ব্যবহার করলে table visual ও আকর্ষণীয় হয়। Hyperlink ব্যবহার করলে table থেকেই অন্য page বা resource-এ যাওয়া যায়। তাই table-এর ভিতরে list, image ও hyperlink ব্যবহার করলে webpage আরও তথ্যবহুল, ব্যবহারবান্ধব ও practical হয়।

\subsection{Common Mistake}

\begin{mistakebox}{সাধারণ ভুল}
\begin{itemize}
\item \texttt{<td>} tag-এর বাইরে list লিখে ফেলা।
\item \texttt{<td>} close করার আগে \texttt{<ol>} বা \texttt{<ul>} ঠিকভাবে close না করা।
\item Image link করতে \texttt{<a>} tag-এর ভিতরে \texttt{<img>} না রাখা।
\item \texttt{href} ও \texttt{src} attribute গুলিয়ে ফেলা।
\item \texttt{<table>} tag-এর ভিতরে সরাসরি \texttt{<ol>} লিখে ফেলা; cell দরকার।
\item Nested table-এর closing tag ভুল করা।
\item Table cell-এর ভিতরে content বেশি দিলে layout অগোছালো করা।
\item External link-এ \texttt{target="\_blank"} লিখতে ভুল করা।
\end{itemize}
\end{mistakebox}

\subsection{১ মিনিটে রিভিশন}

\begin{recapbox}
Table cell-এর ভিতরে list, image, hyperlink ও formatting tag রাখা যায়। List বসাতে \texttt{<td>}-এর ভিতরে \texttt{<ol>} বা \texttt{<ul>} লিখবে। Image বসাতে \texttt{<td>}-এর ভিতরে \texttt{<img>} লিখবে। Link বসাতে \texttt{<td>}-এর ভিতরে \texttt{<a href="">...</a>} লিখবে। Image link করতে \texttt{<a>} tag-এর ভিতরে \texttt{<img>} tag বসাবে। Nested table হলো table-এর cell-এর ভিতরে আরেকটি table।
\end{recapbox}


\begin{boardbox}{পরবর্তী টপিক}
পরের টপিক: Division, Span ও Style Attribute--\texttt{<div>}, \texttt{<span>} এবং inline CSS basic।
\end{boardbox}
